\section{Presentacion convertida en practica: ANCOVA con tres variables explicativas}

Esta seccion nace de una presentacion existente y se incorpora al manual para que el material no permanezca aislado. La presentacion deja de ser un apoyo oral y se transforma en una ruta de estudio: problema, variables, datos, analisis, decision y actividad.

\begin{idea}
Ajustar una respuesta considerando covariables sin confundir control estadistico con causalidad.
\end{idea}

\subsection{Evidencia didactica extraida}

\begin{itemize}
\item Objetivo principal: comprender cómo estas características físicas explican la variabilidad en la aceptación del producto utilizando análisis de covarianza (ANCOVA).
\item La amplitud del rango en la variable de aceptación (17 puntos totales) sugiere una gran variabilidad en las respuestas de los consumidores,
\item lo cual hace especialmente valioso identificar los factores que predicen esta variación.
\end{itemize}

\subsection{Como entra al manual}

El tema se integra al capitulo relacionado con \textbf{ancova}. Si el estudiante solo observa la presentacion, ve una secuencia visual. En el manual, esa misma secuencia se expande con la pregunta de ingenieria, la razon del metodo, la interpretacion y los limites de la conclusion.

\subsection{Practica asociada}

La practica generada pide al estudiante transformar la presentacion en una decision: definir variables, organizar datos, graficar, aplicar el metodo y redactar una conclusion prudente. Si no existia practica, el sistema crea una primera version editable. Si existia una practica relacionada, esta se usa para enriquecer la narrativa y no como bloque aislado.

\subsection{Actividad de evaluacion}

Explique que parte de la presentacion corresponde a contexto, cual corresponde a diseno, cual corresponde a evidencia y cual corresponde a decision. Si falta alguno de estos elementos, proponga como completarlo.
